\documentclass[10pt,conference,a4paper]{IEEEtran}

\usepackage{amsmath,amssymb,mathtools}
\usepackage{graphicx}
\usepackage{microtype}
\usepackage{newtxtext,newtxmath}
\usepackage{float}
\usepackage{balance}
\usepackage{algorithm}
\usepackage{algpseudocode}
\pagestyle{empty}

% ---------------------------------------------------------
% Compact formatting
% ---------------------------------------------------------
\setlength{\parindent}{0pt}
\setlength{\parskip}{1pt}

\setlength{\abovedisplayskip}{3pt}
\setlength{\belowdisplayskip}{3pt}
\setlength{\abovedisplayshortskip}{2pt}
\setlength{\belowdisplayshortskip}{2pt}
\setlength{\jot}{1.5pt}

\setlength{\floatsep}{4pt}
\setlength{\textfloatsep}{4pt}
\setlength{\intextsep}{4pt}
\setlength{\abovecaptionskip}{3pt}
\setlength{\belowcaptionskip}{1pt}

\algrenewcommand\algorithmicindent{1.4em}

% ---------------------------------------------------------
% Commands
% ---------------------------------------------------------
\newcommand{\E}{\mathbb{E}}

\newcommand{\head}[1]{%
    \vspace{2pt}
    \noindent
    \textbf{#1}
    \hspace{2pt}
}

\begin{document}

% ---------------------------------------------------------
% Title
% ---------------------------------------------------------
\title{
Research Problem Formulation
}

\author{
Jad Al Lakkis
}

\maketitle
\thispagestyle{empty}

\vspace{-8pt}

% =========================================================
% PROBLEM
% =========================================================

\head{Problem and main idea.}
A UAV with a $360^{\circ}$ camera cannot send to the user the
live streamed video entirely in high quality due to limited wireless
bandwidth, transmission power, and buffer constraints. So, we
split the video into a base layer and an enhancement layer and,
given a predicted viewport, select specific tiles to enhance using
reinforcement learning. However, this can be expensive, so we use
known system components after the action through post-decision
states to accelerate learning. Then, since the actual viewport and
next channel are statistically independent of an action given the state, virtual experience can be used to further
accelerate learning.

% =========================================================
% SYSTEM FIGURE
% =========================================================

\begin{figure}[H]
    \centering

    \makebox[\columnwidth][c]{%
        \includegraphics[
            width=1.08\columnwidth
        ]{system_model.png}
    }

    \vspace{-4pt}

    \caption{
    Proposed tile-selective aerial $360^{\circ}$ video streaming
    system.
    }

    \label{fig:system}
\end{figure}

% =========================================================
% TILE SELECTION AND ACTION
% =========================================================

\head{Tile selection and action.}
Let the complete $360^{\circ}$ video frame be divided into the
tile set

\begin{equation}
\mathcal{N}
=
\{1,\ldots,N\}.
\end{equation}

The predicted viewport tiles are sent as enhanced tiles. In
addition, the agent selects any other tile from the complete
panorama. The enhanced region is

\begin{equation}
\mathcal{E}^t(\mathbf{x}^t)
=
\left\{
i\in\mathcal{N}:x_i^t=1
\right\}.
\end{equation}

The action is

\begin{equation}
\alpha^t
=
(\mathbf{x}^t,P^t),
\qquad
\mathbf{x}^t
=
(x_1^t,\ldots,x_N^t),
\end{equation}

where $x_i^t=1$ if tile $i$ is enhanced and $x_i^t=0$
otherwise. The predicted viewport tiles are assigned
$x_i^t=1$, while the agent may additionally enhance tiles from
the rest of the panorama. The UAV transmit power satisfies

\begin{equation}
0\leq P^t\leq P_{\max}.
\end{equation}

% =========================================================
% FORCE CMDP TO TOP OF RIGHT COLUMN
% =========================================================

\newpage

% =========================================================
% CMDP
% =========================================================

\head{Constrained Markov decision process.}


\begin{align}
\max_{\pi}\quad
&
\E_{\pi}
\left[
\sum_{t=0}^{\infty}
\lambda^t
\bar Q_{\mathrm{viewport}}^t
\right]
\\[-1pt]
\text{s.t.}\quad
&
\E_{\pi}
\left[
\sum_{t=0}^{\infty}
\lambda^t
\frac{D^t}{T_0}
\right]
\leq
\bar D,
\\[-1pt]
&
\E_{\pi}
\left[
\sum_{t=0}^{\infty}
\lambda^t
\frac{P^t}{P_{\max}}
\right]
\leq
\bar P,
\\[-1pt]
&
\E_{\pi}
\left[
\sum_{t=0}^{\infty}
\lambda^t
\frac{
\sum_{i=1}^{N}x_i^t
}{
N
}
\right]
\leq
\bar B.
\end{align}

Here, $T_0$ is the duration of one time slot,
$\bar Q_{\mathrm{viewport}}^t$ is the normalized viewport quality
at time $t$, and $\bar D$, $\bar P$, and $\bar B$ are the stall,
power, and number-of-tiles budgets, respectively.

The penalty-based reward is

\begin{equation}
r^t
=
\beta_q\bar Q_{\mathrm{viewport}}^t
-
\beta_s\frac{D^t}{T_0}
-
\beta_p\frac{P^t}{P_{\max}}
-
\beta_b
\frac{
\sum_{i=1}^{N}x_i^t
}{
N
}.
\end{equation}

It increases viewport quality while reducing stall time,
transmit power, and unnecessary enhanced tiles.

% =========================================================
% PRE-DECISION STATE
% =========================================================

\head{Pre-decision state.}
The pre-decision state is

\begin{equation}
\omega^t
=
\left(
Z^t,
\Delta_\theta^{t-1},
\Delta_\phi^{t-1},
h^t
\right),
\end{equation}

where $Z^t$ is the playback-buffer level,
$\Delta_\theta^{t-1}$ and $\Delta_\phi^{t-1}$ are the previous
viewport-prediction errors, and $h^t$ is the current
UAV-to-base-station channel gain. The predicted viewport is given seperately using methods like predicting the last viewport where the user was looking , or other possible methods.

% =========================================================
% QOE
% =========================================================

\head{QoE.}
We define the Quality of Experience as follows, ($V_{\mathrm{actu}}^t$ is the actual viewport),
\begin{equation}
C_{\mathrm{cov}}^t
=
\frac{
\left|
V_{\mathrm{actu}}^t
\cap
\mathcal{E}^t(\mathbf{x}^t)
\right|
}{
\left|
V_{\mathrm{actu}}^t
\right|
},
\qquad
\bar Q_{\mathrm{viewport}}^t
=
C_{\mathrm{cov}}^t.
\end{equation}

% =========================================================
% KNOWN QUANTITIES AFTER THE ACTION
% =========================================================

\head{Known quantities after the action.}
Once the action $\alpha^t$ is selected, the enhanced region and
total transmitted data size are known:

\begin{equation}
A^t
=
A_{\mathrm{base}}
+
\sum_{i\in\mathcal{E}^t(\mathbf{x}^t)}
A_{i,\mathrm{enh}}^t.
\end{equation}

The transmission rate is

\begin{equation}
R^t
=
W\log_2
\left(
1+
\frac{P^th^t}{WN_0}
\right).
\end{equation}

The stall time is

\begin{equation}
D^t
=
\begin{cases}
0,
&
Z^t+T_0R^t\geq A^t,
\\[1pt]
T_0
\left(
1-
\dfrac{Z^t+T_0R^t}{A^t}
\right),
&
\text{otherwise}.
\end{cases}
\end{equation}

The deterministic buffer update is

\begin{equation}
\widetilde Z^{t+1}
=
\max
\left\{
Z^t+T_0R^t-A^t,
0
\right\}.
\end{equation}

Therefore, after the action, the following quantities are known:

\[
\mathcal{E}^t(\mathbf{x}^t),
\quad
A^t,
\quad
R^t,
\quad
D^t,
\quad
\widetilde Z^{t+1}.
\]

% =========================================================
% KNOWN / RANDOM SPLIT
% =========================================================

\head{Known/random split and post-decision state.}
The PDS separates the quantities known after the action from the
quantities that are still random:

\[
\omega^t
\xrightarrow{\alpha^t}
\widetilde\omega^t
\xrightarrow{
V_{\mathrm{actu}}^t,\,
h^{t+1}
}
\omega^{t+1}.
\]

The proposed post-decision state is

\begin{equation}
\widetilde\omega^t
=
\left(
\widetilde Z^{t+1},
\Delta_\theta^{t-1},
\Delta_\phi^{t-1},
h^t,
\mathcal{E}^t(\mathbf{x}^t)
\right).
\end{equation}

At the post-decision state, the quantities that remain unknown are

\[
V_{\mathrm{actu}}^t,
\quad
\Delta_\theta^t,
\quad
\Delta_\phi^t,
\quad
\bar Q_{\mathrm{viewport}}^t,
\quad
h^{t+1}.
\]

After observing $V_{\mathrm{actu}}^t$ and $h^{t+1}$, the next
state is

\begin{equation}
\omega^{t+1}
=
\left(
\widetilde Z^{t+1},
\Delta_\theta^t,
\Delta_\phi^t,
h^{t+1}
\right).
\end{equation}

% =========================================================
% REWARD SPLIT
% =========================================================

The reward is then split into a component known immediately after the
action and a component that depends on the random outcome:

\begin{align}
r^t
&=
r_{\mathrm{known}}^t
+
r_{\mathrm{random}}^t,
\\
r_{\mathrm{known}}^t
&=
-\beta_s\frac{D^t}{T_0}
-\beta_p\frac{P^t}{P_{\max}}
-\beta_b
\frac{
\sum_{i=1}^{N}x_i^t
}{
N
},
\\
r_{\mathrm{random}}^t
&=
\beta_q
\bar Q_{\mathrm{viewport}}^t.
\end{align}

% =========================================================
% ROLLOUT / EPISODE STRUCTURE  (NEW, simplified)
% =========================================================

\head{Rollout and episode structure.}
Each episode lasts 36 time steps. A rollout consists of 48
episodes, i.e.\ 1728 environment steps.

% =========================================================
% PSNR  (NEW, with provenance note)
% =========================================================

\head{Viewport PSNR.}
This is the standard PSNR definition (not specific to any of the
references cited below), applied here to this project's own
per-tile, per-quality-level luma MSE data and pooled with the
same viewport-overlap weights already used for
$C_{\mathrm{cov}}^t$. Let $w_i^t\in[0,1]$ denote the fraction of
tile $i$'s area covered by the actual viewport at time $t$, and
$\mathrm{MSE}_i^t$ tile $i$'s luma mean-squared error at whichever
quality level it was actually sent (base or enhanced). The pooled
viewport MSE and PSNR are

\begin{equation}
\overline{\mathrm{MSE}}^t
=
\frac{\sum_i w_i^t\,\mathrm{MSE}_i^t}{\sum_i w_i^t},
\qquad
\mathrm{PSNR}^t
=
10\log_{10}
\!\left(
\frac{255^2}{\overline{\mathrm{MSE}}^t}
\right).
\end{equation}

% =========================================================
% LAGRANGIAN MULTIPLIER UPDATE  (NEW, no more L subscript)
% =========================================================
\head{Lagrangian multiplier update.}
The constraints in the CMDP are expected values, so a Lagrangian
relaxation replaces the fixed weights $\beta_s,\beta_p,\beta_b$ in
$r^t$ with adaptive multipliers $\mu_D,\mu_P,\mu_B\geq0$; from here
on, $r^t$, $r_{\mathrm{known}}^t$, and $r_{\mathrm{random}}^t$ use
these multipliers in place of the fixed weights:

\begin{equation}
r^t
=
\beta_q\bar Q_{\mathrm{viewport}}^t
-\mu_Dc_D^t
-\mu_Pc_P^t
-\mu_Bc_B^t,
\end{equation}

with $c_D^t=D^t/T_0$, $c_P^t=P^t/P_{\max}$, and
$c_B^t=\sum_i x_i^t/N$ (the same three terms as in the CMDP
constraints). For an episode $e$, the discounted constraint
returns are

\begin{equation}
J_j^{(e)}
=
\sum_{t=0}^{35}
\lambda^t c_j^{t,(e)},
\qquad
j\in\{D,P,B\}.
\end{equation}

At the end of each rollout, these returns are averaged over its 48
completed episodes:

\begin{equation}
\overline J_j^{\mathrm{rollout}}
=
\frac{1}{48}
\sum_{e=1}^{48}J_j^{(e)}.
\end{equation}

This sample average provides a Monte Carlo estimate of the expected
discounted constraint return under the current policy. Each
multiplier is then updated once per rollout using the projected
dual-ascent step

\begin{equation}
\mu_j
\leftarrow
\left[
\mu_j
+
\eta_\mu
\left(
\overline J_j^{\mathrm{rollout}}
-
\bar b_j
\right)
\right]_+,
\qquad
j\in\{D,P,B\},
\end{equation}

where $\eta_\mu$ is the multiplier learning rate and
$\bar b_D=\bar D$, $\bar b_P=\bar P$, $\bar b_B=\bar B$.

% =========================================================
% CRITICS  (NEW, no more L subscript)
% =========================================================

\head{Pre-decision and post-decision critics.}
Following the two-critic post-decision-state design of the EHS
paper (adapted here from its soft-actor-critic setting to a
policy-gradient/PPO setting), a critic $V(\omega^t)$ estimates
the pre-decision state and a second critic
$\widetilde V(\widetilde\omega^t)$ estimates the post-decision
state. Since $r_{\mathrm{known}}^t$ is computed analytically
instead of learned from experience, this split leaves each
critic only the part of the value it cannot already compute --
which is what accelerates learning, versus one critic having to
learn the fully blended known-plus-random value from scratch.
Using the known/random split of $r^t$ with the adaptive
multipliers (above):

\begin{equation}
r_{\mathrm{known}}^t
=
-\mu_Dc_D^t-\mu_Pc_P^t-\mu_Bc_B^t,
\qquad
r_{\mathrm{random}}^t
=
\beta_qC_{\mathrm{cov}}^t.
\end{equation}

The two critics are trained toward the targets

\begin{equation}
\widetilde y^t
=
r_{\mathrm{random}}^t+\lambda V(\omega^{t+1}),
\qquad
y^t
=
r_{\mathrm{known}}^t+\widetilde V(\widetilde\omega^t).
\end{equation}

% =========================================================
% ADVANTAGE  (NEW, no more L subscript)
% =========================================================

\head{Advantage.}
Whether the action was good is measured by the one-step
post-decision-state advantage (EHS paper),

\begin{equation}
A^t
=
y^t-V(\omega^t)
=
r_{\mathrm{known}}^t+\widetilde V(\widetilde\omega^t)-V(\omega^t).
\end{equation}

% =========================================================
% ACTOR LOSS / ENTROPY  (NEW, terms labeled)
% =========================================================

\head{Actor loss and entropy regularization.}
The actor is updated with PPO's clipped surrogate objective using
$A^t$ in place of the usual advantage estimate, together with an
entropy regularizer that keeps the policy from collapsing onto a
narrow behavior before the constraints are satisfied:

\begin{equation}
\mathcal L_\theta
=
\underbrace{
-\min\!\big(
\rho^tA^t,\,
\mathrm{clip}(\rho^t,1{-}\epsilon,1{+}\epsilon)A^t
\big)
}_{\text{PPO's own clipped surrogate}}
\;
\underbrace{
-\,\alpha H\big(\pi_\theta(\cdot\,|\,\omega^t)\big)
}_{\text{entropy term, EHS paper}},
\end{equation}

with $\rho^t=\pi_\theta(\alpha^t|\omega^t)/
\pi_{\theta,\mathrm{old}}(\alpha^t|\omega^t)$ and entropy
coefficient $\alpha>0$. The advantage $A^t$ used inside the
clipped surrogate is itself adapted from the EHS paper (previous
section), even though the clipping mechanism around it is PPO's
own.

% =========================================================
% ALGORITHM 1: PDS  (NEW, no more L subscript)
% =========================================================

\begin{algorithm}[H]
\caption{PPO + Post-Decision State Algorithm}
\label{alg:pds}

\small

\begin{algorithmic}[1]
\setlength{\itemsep}{1pt}

\State Initialize actor $\pi_\theta$, critics $V,\widetilde V$,
multipliers $\mu_D,\mu_P,\mu_B$

\Statex \vspace{2pt}

\For{rollout $=1,2,\ldots$}

    \For{$t=0,\ldots,n_{\mathrm{steps}}-1$}

        \State Observe
        $\omega^t=(Z^t,\Delta_\theta^{t-1},
        \Delta_\phi^{t-1},h^t)$

        \State Predict viewport $\widehat V^t$

        \State Sample $\alpha^t\sim\pi_\theta(\cdot\,|\,\omega^t)$

        \State Execute $\alpha^t$ and form the post-decision
        state $\widetilde\omega^t$

        \State Observe $V_{\mathrm{actu}}^t,h^{t+1}$; compute
        $r^t$ and $\omega^{t+1}$

        \State Split $r^t$ into $r_{\mathrm{known}}^t,
        r_{\mathrm{random}}^t$ using $\mu_D,\mu_P,\mu_B$

        \State Store $\big(\omega^t,\alpha^t,\widetilde\omega^t,
        r_{\mathrm{known}}^t,r_{\mathrm{random}}^t\big)$

    \EndFor

    \Statex \vspace{2pt}

    \For{each stored step $t$}

        \State Compute targets $\widetilde y^t,y^t$ and
        advantage $A^t$

    \EndFor

    \State Update $\pi_\theta,V,\widetilde V$ using $A^t$ and
    $\mathcal L_\theta$

    \State Update $\mu_D,\mu_P,\mu_B$ once per rollout

\EndFor

\end{algorithmic}
\end{algorithm}

% =========================================================
% VIRTUAL EXPERIENCE (moved to the end, later work)
% =========================================================

\balance

\head{Virtual experience.}
The actual viewport and next channel are assumed conditionally
independent of a hypothetical action given the current state:

Therefore, the same observed random outcome can be applied to
multiple hypothetical post-decision states:

\[
\left\{
\widetilde\omega^{t,(m)}
\right\}_{m=1}^{M_v}
\xrightarrow{
\substack{
V_{\mathrm{actu}}^t\\
h^{t+1}
}
}
\left\{
\omega^{t+1,m}
\right\}_{m=1}^{M_v}.
\]

Thus, one observed viewport and channel realization can generate
multiple post-decision-state updates. In other words, since
$V_{\mathrm{actu}}^t$ and $h^{t+1}$ do not depend on which
hypothetical action was taken, the single pair we actually
observe is reused, unchanged, to update $\widetilde V$ at every
$\widetilde\omega^{t,(m)}$ in the batch, not only at the
post-decision state of the action that was really executed.

% =========================================================
% ALGORITHM 2: PDS + VIRTUAL EXPERIENCE (moved, unchanged content)
% =========================================================

\begin{algorithm}[H]
\caption{PDS + Virtual Experience Algorithm}
\label{alg:pds-ve}

\small

\begin{algorithmic}[1]
\setlength{\itemsep}{1pt}

\State Initialize the PDS value function $\widetilde V$

\State Choose VE batch size $B_{\mathrm{VE}}$ and VE update period $T$

\Statex \vspace{2pt}

\For{$t=0,1,2,\ldots$}

    \State Observe
    $\omega^t=(Z^t,\Delta_\theta^{t-1},
    \Delta_\phi^{t-1},h^t)$

    \State Predict viewport $\widehat V^t$

    \State Select $\alpha^t=(\mathbf{x}^t,P^t)$ that maximizes
    the known immediate reward plus the estimated
    post-decision-state value.

    \State Execute $\alpha^t$ and form the realized PDS
    $\widetilde\omega^t$

    \State Observe $V_{\mathrm{actu}}^t$ and $h^{t+1}$

    \State Compute coverage, reward $r^t$, and next state
    $\omega^{t+1}$

    \Statex \vspace{2pt}

    \If{$t\bmod T=0$}

        \State Construct the virtual PDS batch

        \Statex
        \[
        \mathcal{B}_{\mathrm{VE}}^t
        =
        \left\{
        \widetilde\omega^{t,(1)},
        \ldots,
        \widetilde\omega^{t,(B_{\mathrm{VE}})}
        \right\}
        \]

        \For{each
        $\widetilde\omega^{t,(b)}
        \in
        \mathcal{B}_{\mathrm{VE}}^t$}

            \State Reuse the observed $V_{\mathrm{actu}}^t$
            and $h^{t+1}$

            \State Compute virtual coverage, reward $r^{t,(b)}$,
            and next state $\omega^{t+1,(b)}$

            \State Update
            $\widetilde V(\widetilde\omega^{t,(b)})$

        \EndFor

    \Else

        \State Update only the realized PDS
        $\widetilde V(\widetilde\omega^t)$

    \EndIf

\EndFor

\end{algorithmic}
\end{algorithm}

\end{document}
